\section{Conclusion}

The four experiments tell a consistent story. Replacing joint with individual taxation lowers the marginal rate on the secondary earner, who responds by working more (Question~3); because human capital accumulates, this response compounds over the life cycle. The extra labor supply broadens the tax base enough that revenue rises by about $48\%$ despite lower rates (Question~4), a Laffer-type result. Holding revenue fixed by recalibrating the tax scale to $\lambda^{\text{indiv},\star} = 1.767$ (Question~5), the household is weakly better off under individual taxation: essentially indifferent when one member starts with no human capital, but clearly better off when both earn --- worth roughly a $3\%$ higher primary wage (Question~6). Individual taxation thus removes a ``marriage penalty'' that is largest precisely when both spouses work.

These conclusions should be read against the model's deliberate simplifications. With no savings there is no consumption smoothing or precautionary motive; the deterministic environment rules out wage risk and the precautionary value of a second earner; and labor supply is purely intensive, with no participation margin or fertility decision. Each of these would tend to strengthen the case for individual taxation by giving the secondary earner additional reasons to work, so the welfare gains here are best read as a conservative lower bound on the mechanism that \cite{borella2023} study in a richer setting.
